\documentclass[a4paper,12pt]{letter}
\usepackage[utf8]{inputenc}
\usepackage{setspace}
\usepackage{geometry}
\usepackage{parskip}   
\usepackage{ragged2e}  
\geometry{left=25mm,right=25mm,top=25mm,bottom=25mm}
\setstretch{1.15}
\signature{\textbf{Zonglong Bai} \\
	On behalf of all co-authors \\
	baizongyao@163.com}
	\address{Zonglong, Bai \\
	North China Electric Power University\\
	Department of Electronic and Communication Engineering \\
	Yonghuabei Street 619th, Baoding, China\\
}
\date{\today}
\begin{document}
\begin{letter}
{Editor-in-Chief \\
	Measurement Science and Technology}
\opening{Dear Editor,}

We would like to submit our manuscript entitled \textit{``AW-DPCNN Based Multi-Representation Acoustic Signal Fusion for Transformer Fault Diagnosis''} for consideration for publication in \textit{Measurement Science and Technology}.

Acoustic sensing has emerged as a promising non-intrusive approach for transformer condition monitoring. However, in practical scenarios, transformer acoustic signals exhibit strong non-stationarity, complex coupling effects, and low signal-to-noise characteristics, which significantly limit the effectiveness of conventional single-representation diagnostic methods. As a result, fully exploiting the information captured by acoustic sensors remains a critical challenge.

In this work, we focus on enhancing the information extraction and representation capability of acoustic sensing data for fault diagnosis. Instead of relying on a single signal perspective, the proposed framework integrates complementary representations to better characterize both time–frequency structures and temporal correlations inherent in transformer acoustic emissions. An adaptive fusion strategy is introduced to improve representation completeness, followed by a deep architecture specifically designed to strengthen multi-scale perception and channel-wise feature selectivity. The main contributions of this work include:
\begin{enumerate}
	\item The development of an adaptive dual-channel fusion mechanism that improves the utilization efficiency of acoustic sensing information by integrating heterogeneous representations in a data-driven manner.
	\item The design of a multi-scale channel-aware deep network that enhances discriminative feature learning, enabling more reliable identification of subtle fault patterns.
	\item Extensive validation on real-world transformer acoustic data, along with independent dataset evaluation on a public benchmark, demonstrates improved diagnostic accuracy, robustness, and generalization capability.
\end{enumerate}

The results suggest that the proposed approach can effectively enhance the reliability of acoustic-based sensing systems for transformer fault diagnosis, providing a practical solution for intelligent condition monitoring in power systems.

This manuscript is original, has not been published previously, and is not under consideration elsewhere. All authors have approved the submission and agree with its contents. We believe that this work fits well within the scope of the Measurement Science and Technology, particularly in the areas of acoustic sensing, intelligent signal processing, and sensor-based fault diagnosis.

Thank you for your time and consideration.

\closing{Sincerely,}

\end{letter}

\end{document}
